\documentclass[12pt]{article}
\usepackage{graphicx} % Required for inserting images
\usepackage{setspace}
\usepackage{fancyhdr}
\usepackage{amsmath}
\usepackage{caption}
\usepackage{graphicx}
\usepackage{booktabs}
\usepackage{adjustbox}
\usepackage{url}
\usepackage[utf8]{inputenc}
\usepackage{csquotes}
\usepackage{natbib}
\usepackage{pdfpages}
\usepackage[raggedright]{ragged2e}
\usepackage[hidelinks]{hyperref}
\usepackage[letterpaper,margin=1in]{geometry}
\bibliographystyle{apsr}   % or plainnat, chicago, apa, etc.
\setcitestyle{authoryear,open={(},close={)},aysep={ }}

\title{Pre-Analysis Plan} 
\author{Wyatt King}
\date{Draft Sections: February 11}

\begin{document}

\maketitle

\section{Theory \& Hypotheses}
As part of my senior thesis, I intend to evaluate three separate but overlapping questions related to public opinion and nuclear security: 1). How do national-identity based arguments restrain public support for nuclear first-use in India, 2). How does knowledge of India's No First Use policy do the same, and 3). Is the combined effect of identity-based arguments and knowledge of India's NFU greater than either effect by its own?

These questions are meant to build on a series of survey experiments that have been conducted over the past ten years, in which scholars have tested how strongly Western public opinion would support the first use of nuclear weapons in different hypothetical scenarios \citep{press_atomic_2013, dill_kettles_2022, smetana_moscow_2022, onderco_hawks_2023, schwartz_when_2024}. Overwhelmingly, these studies have found that the public is sensitive to the military and strategic utility of a nuclear attack in weighing whether or not they approve of it. That being said, there is still reason to believe that a norm against nuclear use exists in some capacity. Moreover, scholars have argued that because public opinion is not particularly sensitive to either international law or the number of civilian casualties incurred by an attack, ideational factors are likely not strongly determinative of public support for nuclear use \citep{sagan_revisiting_2017, sagan_atomic_2024}.


\citet{press_atomic_2013} found, for example, that all else being equal, people vastly prefer the use of conventional weapons to nuclear ones. Moreover, extant studies have largely focused on creating tough tests for the nuclear taboo, in which atomic weapons would achieve a compelling military objective. Yet despite being so useful, nearly half of respondents still oppose nuclear use in most hypotheticals.

\section{Experimental Design}

\section{Test}

To test my hypotheses, I intend to employ a survey experiment in which Indian respondents are presented with a hypothetical crisis between India and Pakistan. The survey will be conducted via Amazon Mechanical Turk because of its cost efficiency relative to other convenience samples. Like other survey platforms in India, MTurk's samples are disproportionately young, male, upper-caste, and wealthy \citep{boas_recruiting_2020}. To assign treatments, I will use a randomized experiment, where respondents are assigned to one of four groups. The control group will receive no information about India's NFU and no arguments condemning nuclear use. One treatment group will receive information about India's NFU but no arguments, and another treatment will read arguments condemning nuclear use but not receive information about India's NFU. Lastly, one treatment group will both be told of India's NFU and read arguments condemning nuclear use.

Drafts of the hypothetical news articles that respondents in each treatment will read are attached in the appendix below. The situation described is consistent with India's Cold Start doctrine and inspired by what security scholars consider to be the most likely case of Indian nuclear first-use \citep{kapur_ten_2008, clary_indias_2019}. Implemented after the Kargil War, Cold Start is Indian military doctrine enabling it to quickly mobilize along the border with Pakistan following a Pakistan-backed terrorist attack. Within a few days, India can then push across the borders with Punjab and Kashmir in an effort either to claim disputed territory of force a favorable compromise in which Pakistan stops sponsoring cross-border terrorism. 

Although India publicly maintains a no first use policy, Christopher Clary and Vipin Narang argue that during such a ground-war India may contemplate launching a first strike nuclear attack. Because India's military is conventionally superior to Pakistan's, the latter may be incentivized to use a small-yield tactical nuclear weapon on the battlefield to hold territory and stop a further Indian incursion. However, given that Pakistan may use a nuclear weapon, India is in turn encouraged to preempt. Fearing retaliation against one of its cities, Clary and Narang argue that India may conduct a counter-force attack, in which it would target Pakistan's nuclear silos, airfields, and missile launchers in an effort to marginalize the chance of a nuclear attack. 

\textit{A priori}, it is difficult to know how successful such an attack would actually be, given uncertainty regarding where Pakistan's nuclear weapons facilities are, how effective its air defenses would be, and a range of other factors. In the article, I note that it is unclear whether all of Pakistan's strategic nuclear weapons have been destroyed, but that Indian government officials speaking on the condition of anonymity believe that there is a low chance of retaliation. I do so as to create a more difficult test for the taboo while not misleading respondents into believing that a counterforce attack would be guaranteed to work. Furthermore, although it is difficult to gauge how many civilians would die in a counterforce attack, I note that the U.N. warns that it is possible hundreds of thousands did. Again, I do so as to not mislead respondents into believing that a counter force attack could be committed with little harm to civilians, even if the extent of that harm is ambiguous.

Within the article, key pieces of information that may affect respondents' support for nuclear use are highlighted in the article title, dek, and pull quotes. Specifically, I emphasize the number of expected casualties, the target of the attack being Pakistan's nuclear weapon facilities, that government officials believe there is little chance of retaliation, and that India launched the attack amid a ground-war with Pakistan.

Across treatment and control groups, the article is identical, except that people assigned to learn about India's no-first use policy read an additional sentence saying that ``the nuclear attack violates long-standing Indian policy that it would never use nuclear weapons unless first attacked with weapons of mass destruction.'' A pull quote in the middle of the article is meant to draw further attention to this fact.

After reading the article, respondents in the control and NFU treatment groups immediately begin answering questions regarding their support for the use of nuclear weapons in the scenario described. People assigned to reading national identity-based arguments, however, are provided three arguments to read criticizing the use of nuclear weapons. Possible arguments that respondents will read are included in the appendix. Broadly, the arguments that I have written are meant to address different aspects of how conceptions of India's national identity impact its foreign policy. Some emphasize India's legacy of non-violence during the nationalist movement, others more explicitly address India's history of advocating for global disarmament, and a third kind emphasize a transcendent identity shared between Indians and Pakistanis. These arguments follow from previous academic literature on how India's strategic culture, national identity, and nuclear weapons policy overlap \citep{hymans_psychology_2006, bajpai_indian_2014}. To decide what survey questions to include in my experiment, I hope to field a short preliminary experiment in which respondents read the vignette and then grade how convincing a randomly assigned article is on a Likert scale. I would then include whatever arguments were graded as the most persuasive in my final experiment. That being said, considering I may lack funding, I may also forgo this trial experiment and use the first three arguments included in the appendix. I would chose these three in particular because they address different aspects of India's history vis-a-vis nuclear weapons and its identity relative to Pakistan.

When it comes to measuring how strongly respondents support the use of nuclear weapons, I intend to follow \citet{sagan_revisiting_2017} and ask respondents two questions: How strongly would they have preferred using nuclear weapons compared to continuing the ground war, and how strongly would they approve of the use of nuclear weapons? Asking respondents the former provides information about what they believe to be reasonable conduct during war. Asking them the latter is a good indication of respondents' willingness to support the prime minister's decision to use nuclear weapons. This level of approval is also politically relevant insofar as respondents should prefer continuing the ground war but still approve of the use of nuclear weapons, then the government likely faces less political risk regardless of public preferences. To further probe the political implications of nuclear use, I intend to ask respondents whether nuclear use would make them more or less likely to vote for the sitting government and whether they would publicly protest the use of nuclear weapons. These questions are meant to be tougher tests for the degree of public outcry to the use of nuclear weapons. 

Note that in keeping with Sagan and Valentino's survey experiments, I intend to use a six-point Likert scale to measure respondents' preference for and approval of nuclear use. On the scale, respondents must either approve or disapprove and cannot indicate that they are unsure or indifferent to the two options. This design mirrors the ``trolley car'' problem, which philosophers use to assess people's moral reasoning when faced with difficult ethical trade-offs \citep{sagan_does_2020}.

In addition to the above, I intend to collect information on potential covariates in a pre-treatment survey administered to respondents. I specifically ask how strongly people approve the Indian National Congress and Bharatiya Janata Party on a seven point Likert scale. These questions will serve as controls for people's political preferences in the regression models described below. I will also ask how strongly respondents agree with the statements that ``India should not use force against another country unless absolutely necessary'' and ``India should never use nuclear weapons unless it is retaliating against another country for using them first.'' Responses to these questions are meant to be highly correlated with their preferences for nuclear use after reading the hypothetical as to reduce standard errors and increase statistical power when I perform my statistical analyses.

To evaluate the support for my hypotheses, I intend to use two statistical procedures primarily. First is a difference of proportions test, comparing people's approval and preference for the nuclear attack between the control and treatment groups. Second is a logistic regression, in which I dichotomize people's approval for nuclear use as the outcome variable and include controls for people's demographics, their political beliefs, and their pre-treatment attitudes toward nuclear use. For robustness, I will also conduct a series of ordinal regressions, using the same independent variables as for the logistic regression.

\section{}

\bibliography{references
}

\end{document}